\documentclass[10pt,a4paper]{amsart}
\usepackage[margin=19mm]{geometry}
\usepackage{amsmath,amssymb,mathtools}
\usepackage[scr=rsfs,cal=euler]{mathalfa}
\usepackage{xfrac}
\usepackage[unicode,hidelinks,bookmarks=false]{hyperref}
\newcommand{\ie}{i.e.\ }
\newcommand{\cf}{cf.\ }
\newcommand{\resp}{respectively\ }
\newcommand{\Subsection}{\subsection}
\providecommand{\nobreakdash}{\nobreak\hbox{-}}
\setlength{\emergencystretch}{2em}
\multlinegap0pt
\usepackage{amssymb}
\usepackage{tikz}
\usepackage{float}
\usepackage{mathtools}
\usepackage{dsfont}
\newtheorem{theorem}{Theorem}
\newtheorem{claim}{Claim}
\DeclareMathOperator{\supp}{supp}
\title{Explicit constructions of optimal blocking sets and minimal codes}
\author{Anurag Bishnoi, Istv\'an Tomon}
\date{Source: arXiv:2411.10179v2 (CC BY 4.0)}
\begin{document}
\maketitle

\noindent\emph{Source and licence.} Explicit constructions of optimal blocking sets and minimal codes. Version arXiv:2411.10179v2 (CC BY 4.0), distributed by arXiv under the Creative Commons Attribution 4.0 International licence, http://creativecommons.org/licenses/by/4.0/, as recorded in the arXiv licence metadata for this exact version and shown on the abstract page https://arxiv.org/abs/2411.10179v2. Full licence notice: \texttt{licenses/arxiv-2411.10179-CC-BY-4.0.txt}.\par\smallskip

\noindent\textit{Editorial scope.} Counted unit: the article's theorem theorem (source0.tex lines 954--957). The article's complete original proof of it is delivered with it (source0.tex lines 959--1009, one \texttt{proof} environment of 51 lines). No mathematics is written, repaired or re-derived: the statement and the proof are the article's own lines, in the article's own notation, and no retained line is substituted or re-flowed.  No further source-local context is needed: every cross-reference of every retained line resolves inside the two delivered spans.  Nothing was omitted from any retained span, so the package carries no omission record.  No retained line contains a citation, so the wrapper carries no bibliography and no bibliography entry.  No retained line contains a figure environment, an image-inclusion command or drawing code, so no image is delivered and the package declares no asset dependency.  Editorial formatting only, in addition: an AMS \texttt{amsart} preamble that restates the article's own declarations and macros used by the retained lines, namely the environments \texttt{theorem}, \texttt{claim} on separate counters, together with the standard packages those lines need.  The retained environments print in this document as Theorem 1, Claim 1 in order of appearance, whereas the article prints the same items with the numbering of its own full document.  The complete native source of the article as distributed in the arXiv e-print for this exact version is bundled unchanged as \texttt{sources/168588/source0.tex}.
\begin{theorem}
    For every integer $s$ and $\varepsilon>0$, there exists $C$ such that the following holds. 
Let $q$ be a prime power and $k$ be an integer such that $k\geq \varepsilon q$, and both are sufficiently large with respect to $s$ and $\varepsilon$. Then there exists a set $W\subset \mathbb{F}_q^k$ of size at most $Ck$ such that any $s+1$ elements of $W$ are linearly independent, and any $\varepsilon |W|$ elements of $W$ span $\mathbb{F}_q^k$. Moreover, such a set $W$ can be constructed in $k^{O_{s,\varepsilon}(1)}$ time.
\end{theorem}

\begin{proof}
Choose  $R \in (0, 1)$ and an integer $M \ge 2$ strictly depending on $\varepsilon$, such that $R+ \frac{1}{M}< \frac{\varepsilon}{2}$.
Let $m$ be the minimal integer such that $m \ge s+2$ and $q^m \ge \frac{k}{mR}$. Define $Q = q^m$, and let $K = k/m$ (where we assume that $m$ divides $k$, the general case being almost identical, but slightly more technical), and  $N = \lceil K/R \rceil$. By our choice of $m$, we have $Q>N$, so we can choose a set $S = \{\alpha_1, \dots, \alpha_N\}$  of $N$ distinct elements of $\mathbb{F}_Q^*$.

Fix an arbitrary basis of $\mathbb{F}_Q$ over $\mathbb{F}_q$ and define the canonical linear isomorphism $\phi : \mathbb{F}_Q \rightarrow \mathbb{F}_q^m$.

\begin{claim}
There exists a linear mapping $L : \mathbb{F}_Q \to \mathbb{F}_{q}^{mM}$ of the form $L(x) = \left(\phi(\gamma_1 x), \dots, \phi(\gamma_M x)\right)$ such that for all $x \in \mathbb{F}_Q^*$, the  fraction of non-zero coordinates of $L(x)$ is strictly bounded below by $\delta_{0} = 1 - 2/M$. Moreover, if $v\in \mathbb{F}_q^{mM}\setminus\{0\}$ with at most  $s+1$ non-zero entries, then $\langle v,L(x)\rangle\neq 0$ for some $x\in \mathbb{F}_Q^*$.
\end{claim}
\begin{proof} Pick $\vec{\gamma} = (\gamma_1, \dots, \gamma_M) \in (\mathbb{F}_Q)^M$ randomly from the uniform distribution, and define $L_{\vec{\gamma}}$ accordingly. Let $d = \lfloor \delta_{0} mM \rfloor$. For any fixed non-zero $x \in \mathbb{F}_Q^*$,  $L_{\vec{\gamma}}(x)$ is uniformly distributed over $(\mathbb{F}_q^{m})^M$. The probability that $L_{\vec{\gamma}}(x)$ has less than $d$ non-zero entries is the volume of the Hamming ball $V_q(mM, d-1)$ divided by $q^{mM}$. We apply the union bound over all $Q-1$ non-zero elements of $\mathbb{F}_Q^*$
$$\mathbb{P} \left[ \exists x \neq 0 : |\supp(L_{\vec{\gamma}}(x))|<d\right] \le (q^m - 1) \frac{V_q(mM, d-1)}{q^{mM}} < q^m \frac{q^{mMH_q(\delta_{0})}}{q^{mM}},$$
where $$H_q(t) = t \log_q(q-1) - t \log_q(t) - (1-t) \log_q(1-t)$$ is the $q$-ary entropy function. As $\lim_{q\rightarrow\infty} H_q(t)=t$, in case $\delta_0=1-2/M$, we have $H_q(\delta_0)< 1-1.5/M$ for $q$ sufficiently large, so the right-hand-side is less than 1/2.

The number of vector $v$ with at most $s+1$ non-zero entries is at most $q^{s+1}(mM)^{s+1}$. The probability that for a fixed non-zero $v$, we have $\langle v,L_{\vec{\gamma}}(x)\rangle=0$ for every $x$ is $1/q^m$. Therefore, by the union bound
$$\mathbb{P} \left[ \exists v \neq 0 : |\supp(v)|\leq s+1\text{ and }\forall x: \langle v,L_{\vec{\gamma}}(x)\rangle=0\right] \le \frac{q^{s+1}(mM)^{s+1}}{q^m}.$$
Thus, if $m>s+1$ and $q$ is sufficiently large, the right-hand-side is less than $1/2$.
Hence, at least one valid  $\vec{\gamma}$ exists.
\end{proof}

As $Q$ is the smallest power of $q$ larger than $\max\{q^{s+2},k/(mR)\}$, we have $Q\leq k^{2s}$ for $k\geq \varepsilon q$ sufficiently large. Therefore, by searching all $M$-tuples $(\gamma_1,\dots,\gamma_M)$ of $\mathbb{F}_Q$, we can find the suitable linear mapping $L$ in $k^{O(sM)}$ time.


Let $\mathcal{P}_K$ be the $\mathbb{F}_Q$-vector space of polynomials with degree strictly less than $K$, the $\mathcal{P}_K\cong \mathbb{F}_q^{Km}=\mathbb{F}_q^k$. Define the $\mathbb{F}_q$-linear transformation $\Phi : \mathcal{P}_K \to \mathbb{F}_q^{NmM}$ as $$\Phi(P) = \left( L(P(\alpha_1)), L(P(\alpha_2)), \dots, L(P(\alpha_N)) \right).$$
Let $n = NmM$. By fixing an $\mathbb{F}_q$-basis for $\mathcal{P}_K$, the transformation $\Phi$ is represented by a matrix $A \in \mathbb{F}_q^{n \times k}$, such that $\Phi(y) = Ay$ for any coefficient vector $y \in \mathbb{F}_q^k$. We define $W \subset \mathbb{F}_q^k$ as the set of the $n$ columns of the transposed matrix $A^T$.

\begin{claim}
$|W| \le Ck$ for a constant $C$ depending only on $\varepsilon$.
\end{claim}

\begin{proof} We have $|W| = n = NmM$ and $k=mK$. Therefore, $|W|/k=NM/K\leq (K/R+1)M/K=M/R+M/K<M/R+1$. As $R$ and $M$ only depend on $\varepsilon$, this proves the claim.
\end{proof}

\begin{claim}
 If $k$ is sufficiently large with respect to $s,\varepsilon$, then any $s+1$ elements of $W$ are linearly independent over $\mathbb{F}_q$.
 \end{claim}
 
 \begin{proof} Assume to the contrary that this is false, then there exists a non-zero vector $v \in \mathbb{F}_q^n$ with at most $s+1$ non-zero entries such that $A^T v = 0$. This implies that for every $P\in \mathcal{P}_K$, we have $\langle v,\Phi(P)\rangle=0$. Recall that $\Phi(P)=(L(P(\alpha_1)),\dots,L(P(\alpha_N))$. Write $v=(v_1,\dots,v_N)$, then 
$\langle v,\Phi(P)\rangle=\sum_{i=1}^N \langle v_i,L(P(\alpha_i))\rangle$.
Let $j\in [N]$ such that $v_j\neq 0$. As $|\supp(v_j)|\leq s+1$, there exists some $x_j$ such that $\langle v_j,L(x_j)\rangle\neq 0$. Also, there are at most $s+1$ indices $i\in [N]$ such that $v_i\neq 0$. Therefore, we can find a polynomial of degree at most $s < K$ (which is satisfied if  $k$ is sufficiently large with respect to $s,\varepsilon$) such that $P(\alpha_j)=x_j$ and $P(\alpha_i)=0$ for all indices $i\neq j$ for which $v_i \neq 0$. But then $\langle v,\Phi(P)\rangle=\langle v_j,L(x_j)\rangle\neq 0$.
 \end{proof}
 
 
\begin{claim} 
Any subset of $\varepsilon|W|$ elements of $W$ spans $\mathbb{F}_q^k$
\end{claim} 
\begin{proof} 
If $U\subset W$ does not span $\mathbb{F}_q^k$, then there exists a non-zero $v\in \mathbb{F}_q^k$ orthogonal to all elements of $U$. Equivalently, $Av(u)=0$ for every $u\in U$. But observe that we can write  $Av=\Phi(P)$ for some $P\in \mathcal{P}_K$. Here, $\Phi(P)=(L(P(\alpha_1)),\dots,L(P(\alpha_N))$. As $P$ has degree less than $K$, we have $P(\alpha_i)=0$ for at most $K-1$ different indices $i\in [N]$. Moreover, if $P(\alpha_i)\neq 0$, then $L(P(\alpha_i))$ has at least $\delta_0$ non-zero coordinates. In total, we have
$$|U|\leq (K-1)mM+(1-\delta_0)n=\left(\frac{K-1}{N}+\frac{2}{M}\right)n\leq \left(2R+\frac{2}{M}\right)n\leq \varepsilon n.$$
\end{proof}

\end{proof}

\end{document}
